\documentclass[a4paper,12pt]{article}
\usepackage[utf8]{inputenc}
\usepackage[T1]{fontenc}
\usepackage[english]{babel}
\usepackage{geometry}
\geometry{left=2.5cm,right=2.5cm,top=2.5cm,bottom=2.5cm}

\usepackage{lmodern}
\usepackage{xcolor}
\definecolor{titleblue}{RGB}{0,51,140}

\usepackage{hyperref}

\pagestyle{empty}

\begin{document}

\begin{flushleft}
    \textbf{Ismail AISSAOUI} \\
    State Engineer in Artificial Intelligence \\
    ismail.aissaoui.pro@gmail.com \\
    +213 660 70 77 96 \\
    \href{https://ismail-aissaoui.vercel.app}{ismail-aissaoui.vercel.app}
\end{flushleft}

\begin{flushright}
    To the PhD Selection Committee \\
    \textbf{EDT Program - IMT Atlantique / IRISA} \\
    Attn: Prof. Antoine Beugnard and Prof. Jean-Marc Jézéquel \\
    \medskip
    May 1, 2026
\end{flushright}

\bigskip

\noindent \textbf{Subject: Application for PhD: Federation of data and models to define digital twins (Rank 6 - PC2)}

\bigskip

Dear Prof. Beugnard and Prof. Jézéquel,

As a State Engineer in Artificial Intelligence from ESI-SBA, I am writing to express my enthusiastic application for the PhD position in "Federation of data and models to define digital twins" within the Engineering Digital Twin (EDT) national program. This application is motivated by the fit between my expertise in knowledge extraction and heterogeneous data modeling and your research on long-term digital twin repositories.

Currently, I am completing my thesis at Sonatrach, where I have developed a **Physics-Informed Digital Twin** for gas turbine monitoring. A core challenge of my project has been the integration of highly heterogeneous information—ranging from high-frequency sensor streams and 3D CAD models to behavioral physics equations and historical maintenance logs. This experience has given me a deep understanding of the **semantic and temporal heterogeneity** that currently hampers the creation of unified digital twin knowledge bases.

I am particularly excited by your proposal to develop an **information federation approach** to structure these heterogeneous sources. My technical background in **Semantic Knowledge Retrieval (RAG)** and high-dimensional statistical analysis aligns perfectly with the goal of identifying and classifying descriptive, prescriptive, and predictive models. I am eager to contribute to the development of a unified vocabulary (ontology) that enables navigation across the multidimensional space of time and variants, ensuring that digital twins remain maintainable and scalable throughout their entire lifecycle.

My background in semantic retrieval and my experience with production-grade PyTorch and Model-Based Systems Engineering (MBSE) provide me with the necessary tools to contribute to the **Artemis platform** and the broader EDT ecosystem. I am a highly autonomous researcher, comfortable bridging the gap between software engineering principles and industrial digital twin repositories.

Thank you for your time and consideration. I look forward to the possibility of discussing how my experience in data federation and hybrid modeling can contribute to the success of your research group in Brest.

Sincerely,

\bigskip

\textbf{Ismail AISSAOUI}

\end{document}
